% Standalone problem document, generated from the adjacent README.md.
% Compile with XeLaTeX. Mathematical content is maintained in Markdown.
\documentclass[11pt,a4paper]{article}
\usepackage[fontsize=10.5pt]{scrextend}
\usepackage[margin=22mm,headheight=15pt,headsep=9mm,footskip=11mm]{geometry}
\usepackage{fontspec}
\setmainfont{texgyrepagella}[Extension=.otf,UprightFont=*-regular,BoldFont=*-bold,ItalicFont=*-italic,BoldItalicFont=*-bolditalic]
\setsansfont{texgyreheros}[Extension=.otf,UprightFont=*-regular,BoldFont=*-bold,ItalicFont=*-italic,BoldItalicFont=*-bolditalic]
\setmonofont{texgyrecursor}[Extension=.otf,UprightFont=*-regular,BoldFont=*-bold,ItalicFont=*-italic,BoldItalicFont=*-bolditalic,Scale=MatchLowercase]
\usepackage{amsmath,amssymb}
\usepackage{unicode-math}
\setmathfont{texgyrepagella-math.otf}
\usepackage{xcolor,microtype,enumitem,needspace,fancyhdr}
\usepackage{bookmark}
\definecolor{ink}{HTML}{17344B}
\definecolor{accent}{HTML}{197D80}
\definecolor{muted}{HTML}{596773}
\hypersetup{colorlinks=true,linkcolor=accent,urlcolor=accent,pdftitle={TR-10 - Border
Comon's conjecture over the complex
numbers},pdfauthor={Open Problems in Numerical Linear Algebra}}
\urlstyle{same}
\setlength{\parindent}{0pt}
\setlength{\parskip}{5pt plus 1pt minus 1pt}
\linespread{1.04}
\setlength{\emergencystretch}{3em}
\widowpenalty=10000
\clubpenalty=10000
\displaywidowpenalty=10000
\allowdisplaybreaks[1]
\setlist{nosep,leftmargin=1.5em}
\providecommand{\tightlist}{\setlength{\itemsep}{0pt}\setlength{\parskip}{0pt}}
\providecommand{\pandocbounded}[1]{#1}
\pagestyle{fancy}
\fancyhf{}
\fancyhead[L]{\sffamily\footnotesize\color{muted}OPEN PROBLEMS IN NUMERICAL LINEAR ALGEBRA}
\fancyhead[R]{\sffamily\bfseries\footnotesize\color{accent}TR-10}
\fancyfoot[L]{\parbox[b]{0.9\textwidth}{\sffamily\fontsize{8}{10}\selectfont\color{muted}Editorial ratings; a bounded search cannot certify that a problem remains unsolved.\\Literature check: 2026-09-10}}
\fancyfoot[R]{\sffamily\footnotesize\color{muted}\thepage}
\renewcommand{\headrulewidth}{0.3pt}
\renewcommand{\headrule}{\hbox to\headwidth{\color{accent}\leaders\hrule height \headrulewidth\hfill}}
\makeatletter
\renewcommand{\section}{\Needspace{5\baselineskip}\@startsection{section}{1}{0pt}{12pt plus 2pt minus 2pt}{4pt}{\sffamily\bfseries\large\color{ink}}}
\renewcommand{\subsection}{\Needspace{4\baselineskip}\@startsection{subsection}{2}{0pt}{9pt plus 1pt minus 1pt}{3pt}{\sffamily\bfseries\normalsize\color{ink}}}
\makeatother
\setcounter{secnumdepth}{0}
\begin{document}
{\sffamily\small\color{muted}Tensor computations\par}
\vspace{3pt}
{\raggedright\hyphenpenalty=10000\exhyphenpenalty=10000\sffamily\bfseries\fontsize{21}{24}\selectfont\color{ink}Border
Comon's conjecture over the complex numbers\par}
\vspace{7pt}
{\sffamily\small\textbf{Difficulty:} extreme\quad\textbf{Importance:} interesting
to the community\par}
{\sffamily\small\bfseries\color{accent}Status: Partially resolved\par}
\vspace{4pt}
{\color{accent}\hrule height 0.8pt}
\vspace{6pt}
\textbf{Rating rationale:} Extreme because equality of two secant-rank
notions for all symmetric tensors remains a foundational tensor-geometry
barrier; community importance concerns whether symmetry can be imposed
without an approximation-rank penalty.

\subsection{Statement}\label{statement}

For every integer order \(d\ge3\), dimension \(n\ge2\), and symmetric
tensor \(T\in\operatorname{Sym}^d(\mathbb C^n)\), is \[
\underline R(T)=\underline R_{\mathrm{sym}}(T)?
\] Here \(R(T)\) is the least number of summands in a decomposition
\(T=\sum_{j=1}^r v_{j,1}\otimes\cdots\otimes v_{j,d}\), with complex
vectors. Symmetric rank restricts each summand to
\(c_jv_j^{\otimes d}\). The border rank \(\underline R(T)\) is the least
\(r\) for which a sequence of tensors of rank at most \(r\) converges to
\(T\) in the entrywise Euclidean topology. Symmetric border rank instead
uses symmetric tensors of symmetric rank at most \(r\).

\subsection{Relevance}\label{relevance}

This asks whether exploiting symmetry can increase the number of
rank-one terms required for arbitrarily accurate approximation. It
concerns the attainable accuracy of symmetric low-rank tensor models.

\subsection{References}\label{references}

\begin{enumerate}
\def\labelenumi{\arabic{enumi}.}
\tightlist
\item
  T. Mańdziuk and E. Ventura, \emph{Symmetrization maps and minimal
  border rank Comon's conjecture}, arXiv:2411.05721v2, 2026.
  \href{https://arxiv.org/html/2411.05721v2}{Full text}, §1, opening
  discussion and Conjecture 1.1; the general conjecture precedes the
  minimal-border-rank specialization.
  \href{https://arxiv.org/abs/2411.05721}{Version record}.
\item
  J. M. Landsberg, \emph{Geometry and Complexity Theory}, Cambridge
  University Press, 2017.
  \href{https://people.tamu.edu/~jml/simonsclass.pdf}{Author text},
  Conjecture 5.6.1.5, a three-factor formulation.
\end{enumerate}

\subsection{Status check --- 2026-09-10}\label{status-check-2026-09-10}

Rechecked \href{https://arxiv.org/html/2411.05721v2}{Mańdziuk--Ventura
v2, §1 and the main theorems}, and searched for 2025--2026 proofs or
counterexamples. The September 2026 revision explicitly leaves the
general border conjecture open while proving minimal-border-rank
subfamilies. Those are genuine portions of this universal target;
Shitov's counterexample to ordinary exact-rank Comon is a different
statement. No full border-rank resolution was located.
\end{document}
